% !TEX root = main.tex
\section{Schedule and Management Plan}


\begin{newitemize}
  \item \textbf{Year 1 – Tight threshold Schnorr.} 
  Formalize the Vandermonde Circular DL assumption (VCDL) and relation to CDL/LDVR; 
  prove security for $\mathsf{Sparkle+}$ and $\mathsf{Gorgos}$ threshold Schnorr; impossibility for $\mathsf{FROST}$; 
  lay groundwork for post-quantum instantiability results.
%  \textbf{Deliverables:} (i) definitions/relations, (ii) proofs, (iii) ESRP integration. 

  \item \textbf{Year 2 – Core instantiability work.} 
  Develop tunable lossy functions (TLFs) and adaptive-leakage UCEs including definitions and lattice-based constructions, and obtain standard-model instantiations of $\mathsf{FO}$ and variants 
  across all NIST PQC PKE schemes. 
 % \textbf{Deliverables:} (i) definitions, (ii) instantiations, 
%  (iii) materials forming a graduate seminar. 

  \item \textbf{Year 3 – Extending  tightness.} 
  Generalize CDL to lattices via circular-MSIS for analysis of $\mathsf{Dilithium}$; 
  design CDL variants for analysis of $\mathsf{Bulletproofs}$-based range proofs. 
 % \textbf{Deliverables:} (i) definitions/relations, 
%  (ii) $\mathsf{Dilithium}$ proofs, (iii) proofs for $\mathsf{Bulletproofs}$-based range proofs. 
\end{newitemize}

%The PI will (a) supervise all researchers, 
%(b) integrate results annually  into  core cryptography courses and a Year-2 graduate seminar, 
%and (c) coordinate with outreach programs.
%\noindent In tandem, the PI will teach two cryptography courses (undergraduate and graduate) each year, embedding our  results as appropriate. Each phase aligns also with educational and outreach deliverables, ensuring research outputs directly train students. %The PI will meet weekly with students and assess progress against milestones.
%\heading{Standards touchpoints and artifacts.}
%Year~1: memo to NIST Threshold WG with tight static/adaptive tradeoffs and parameter guidance for Sparkle$^+$ and Gorgos. 
%Year~2: practice-oriented note on FO instantiation assumptions for PQC KEMs (alignment with SHA-2/3 properties). 
%Year~3: circular-MSIS analysis for Dilithium’s ROM tightness. 
%All results will ship with ePrint preprints, proof scripts, and test vectors in an open repository.

 %\textbf{Dissemination and collaboration.} Each year’s research outputs will be shared promptly via ePrint and conference submissions (Crypto, Eurocrypt, TCC, PKC), and discussed with NIST collaborators (e.g., Donna Dodson, John Kelsey) to inform the ongoing threshold-signature and PQC standardization efforts. Outreach activities (UMass Turing Summer Program, ESRP, AISec, and the Massenberg Summer STEM Institute) will be aligned with research milestones, ensuring timely transfer of ideas into education and community engagement.

%\textbf{Teaching integration.} As PI I will teach one graduate and one undergraduate cryptography course each year, with a new advanced topics seminar in Year~2 focused on instantiability and threshold cryptography.


%\noindent This structure ensures steady research progress, regular deliverables to the cryptography community, and sustained educational and societal impact throughout the award period.

 \iffalse

\noindent \textbf{Year One:} \underline{Research:} Sections 2.1, 2.2, 3.1, 4.1, 4.2. \underline{Teaching:} Fall --- Graduate introduction to cryptography.   Spring --- Undergraduate introduction to cryptography. \\
\textbf{Year Two:}
\underline{Research:}   Sections 2.3, 3.2, 3.3, 4.3.  \underline{Teaching:} Fall --- Graduate introduction to cryptography.  Spring --- Undergraduate introduction to cryptography, advanced topics seminar. \\
\textbf{Year Three:} \underline{Research:} Sections  2.4, 3.2, 3.3, 4.4.  \underline{Teaching:} Fall --- Graduate introduction to cryptography.  Spring --- Undergraduate Introduction to cryptography. 

\fi 
\newpage
\setcounter{page}{1}

%Additionally, the PI is planning to apply for tenure in year two, which if approved would introduce an additional course in subsequent years (as the department's current teaching load for tenured faculty is 2-1), probably a seminar course.
%As seen above, the PI plans to roughly proceed through the components in order, 
%e has   allocated the most time for Component I, which is expected to be the most involved.
%As the components do not really have interdependencies, he will  maintain some flexibility as to their ordering to accommodate student interest and collaborator availabilities.   %ßß?ç
%We will maintain  flexibility in the research schedule according to availability of collaborators, interests of students, \emph{etc.}
%Component I-C is well-suited for student with a strong combinatorics and algorithms background, while most other Components will require a strong cryptography background.    We also hope to allocate a post-doc to Components I or III.
%Undergraduate  and high school research  will focus on the part of Component III-B that looks at program obfuscation techniques in practice.
%In undertaking this ambitious research proposal, the PI will continue his collaborations with other major research  groups in the cryptographic  community.  A 
%The PI plans also to put  focus on collaborating with  ``applied'' researchers (\emph{e.g.}, George Kollios, Dave Levin, Nick Feamster,  Micah Sherr) as well as  the Georgetown University Medical Center and NIST,  not only to  guide our initial  goals but also integrate the results into real systems.